\documentclass[answers]{exam}

% Version compilable
%\usepackage{../../../mypackages}
%\usepackage{../../../macros}

% Version pour execution du script
\usepackage{../../mypackages}
\usepackage{../../macros}

% Pour le tableau
\usepackage{array}
\usepackage{longtable}
\usepackage{multirow}
\usepackage[table]{xcolor}

\title{Devoir en classe - Geogebra + Tracés géométriques}
\author{Nicolas Bancel}
\date{Mardi 9 Décembre 2025}

% Mise en forme des solutions en bleu
\SolutionEmphasis{\color{blue}}
\renewcommand{\solutiontitle}{\noindent}

\definecolor{lavande}{RGB}{180,190,255}

% Colonne p{} alignée à gauche et qui wrappe
\newcolumntype{L}[1]{>{\raggedright\arraybackslash}p{#1}}

\begin{document}

\textbf{Collège Lycée Suger}
\hfill
\textbf{Mathématiques} \\

\textbf{Année 2025-2026}
\hfill
\textbf{Classe de 6ème} \par

{\let\newpage\relax\maketitle}

\begin{tcolorbox}[
  %colback=blue!5!white,
  colframe=blue!90!black,
  %title=Commentaire plus détaillé,
  fonttitle=\bfseries,
  boxsep=4pt,
  sharp corners,
]
\begin{center}
\textbf{\textcolor{blue}{Durée de l'exercice : 1h40 \\
Coefficient : 0.5 \\
Note sur 20 points
}}
\end{center}
\end{tcolorbox}


\section*{Angles dans le triangle — La propriété des 180°}

\begin{questions}

\question[2] Sur Geogebra : 
\begin{parts}
  \part Tracer un triangle et mesurer ses trois angles.
  \part Faire la somme de ces mesures.
  \part Comparer ce résultat avec les autres élèves de la classe. Que constate-t-on ?
  \part Exporter le fichier Geogebra en le sauvegardant sur votre tablette au format .ggb, et en le nommant \textit{Prénom Nom - Triangle avec Angle}
\end{parts}

\question[2] Sur une feuille de papier (quadrillée)
\begin{parts}
  \part Tracer un nouveau triangle "assez grand" et le découper avec précision en suivant son contour.
  \part Nommer ses trois angles $\hat{A}$, $\hat{B}$ et $\hat{C}$.

\begin{figure}[H]
  \centering
  \includegraphics[width=0.5\textwidth]{activite_1_01.jpg}
  \caption{Exercice Geogebra}
\end{figure}

  \part Réaliser un pliage afin de rabattre les trois sommets du triangle vers l'intérieur pour former un rectangle.

  \begin{figure}[H]
  \centering
  \includegraphics[width=0.5\textwidth]{activite_1_02.jpg}
  \caption{Exercice Geogebra}
\end{figure}

  \part Expliquer pourquoi ce pliage permet d'expliquer le résultat établi dans la partie 1.
\end{parts}

\end{questions}

\section*{A la recherche du cercle magique}

\begin{questions}

\question \textbf{Partie 1} \\
Sur papier, reproduire le triangle ci-dessous et essayer à l'aide du compas de construire un cercle qui passe par les trois sommets du triangle.


  \begin{figure}[H]
  \centering
  \includegraphics[width=0.5\textwidth]{activite_2_01.jpg}
  \caption{Exercice Geogebra}
\end{figure}

\question \textbf{Partie 2}
\begin{parts}
  \part Sur papier, reproduire à nouveau le triangle ci-dessus.
  \part Construire les médiatrices des trois côtés du triangle.
  \part Que constate-t-on ?
\end{parts}

\question \textbf{Partie 3}  
Les médiatrices se coupent en un point appelé point de concours.  
On dit que les trois médiatrices sont concourantes.
\begin{parts}
  \part Nommer O le point de concours.
  \part Tracer le cercle de centre O et passant par A.
  \part Que constate-t-on ?
\end{parts}

\question \textbf{Partie 4}  
Le cercle tracé à la partie 3 s'appelle le cercle circonscrit au triangle ABC.
\begin{parts}
  \part Tracer un triangle équilatéral et tracer son cercle circonscrit. Que constate-t-on ?
  \part Faire de même avec un triangle rectangle. Que constate-t-on ?
\end{parts}

\end{questions}


\section*{Rendu du devoir}

\begin{questions}

\question Aller sur EcoleDirecte, et envoyer un email dont le titre est "Prénom Nom - Devoir Geogebra" et qui contient en pièce jointe le fichier GeoGebra de l'exercice 1 (Angles dans le triangle)
\question Puis coller sur une copie double les différents dessins et commentaires faits dans les 2 exercices. Vous inscrirez votre prénom et votre nom sur la copie double.
\end{questions}

\section*{Barème}
\renewcommand{\arraystretch}{1.7}

{
  %\setlength{\LTleft}{-2.3cm}  % déborde de 1.5 cm dans la marge de gauche
  \setlength{\LTright}{0pt}    % rien de plus à droite

\centering
\begin{longtable}{|L{0.3\textwidth}|L{0.7\textwidth}|c|}
\hline
\rowcolor{lavande}
\textbf{Critère évalué} & \textbf{Description} & \textbf{Points} \\[2pt]
\hline

% ---------------- Exercice 1 ----------------
  Exercice 1 — Angles dans le triangle (GeoGebra + feuille)
  & Construction correcte d'un triangle sur GeoGebra, mesure des trois angles et calcul de la somme.  
  Fichier GeoGebra enregistré et nommé correctement ("Prénom Nom - Triangle avec Angle").  
  Sur feuille : triangle tracé proprement, pliage réalisé, explication cohérente du lien avec la somme des angles.
  & 5 \\
\hline

% ---------------- Exercice 2 ----------------
  Exercice 2 — "À la recherche du cercle magique"
  & Reproduction correcte des triangles demandés.  
  Médiatrices tracées avec soin, point de concours identifié et nommé O.  
  Cercle de centre O passant par A correctement tracé.  
  Cercles circonscrits des triangles équilatéral et rectangle réalisés, avec constats rédigés clairement pour chaque cas.
  & 5 \\
\hline

% ---------------- Propreté ----------------
  Propreté et soin du travail
  & Écriture lisible et assez régulière.  
  Figures tracées aux instruments (règle, équerre, compas), pas de tracés à main levée.  
  Peu de ratures, feuilles propres, découpages et collages soignés.
  & 3 \\
\hline

% ---------------- Rendu ----------------
  Rendu du devoir (email + copie double)
  & Email envoyé via ÉcoleDirecte avec un objet clair du type "Prénom Nom - Devoir Geogebra".  
  Fichier GeoGebra joint au bon format (.ggb).  
  Copie double rendue avec prénom, nom, classe, et tous les dessins / réponses collés ou recopiés proprement.
  & 3 \\
\hline

% ---------------- Entraide et questions ----------------
  Entraide et qualité des questions posées
  & L'élève demande de l'aide de manière précise (par exemple : "Je n'arrive pas à construire les médiatrices" plutôt que "Je n'ai pas compris").  
  Il ou elle peut aider un camarade en expliquant une méthode sans simplement donner la réponse.  
  Les questions posées montrent une vraie recherche de compréhension.
  & 4 \\
\hline

% ---------------- Total ----------------
\rowcolor{green!35}
 \textbf{Total} & \textbf{Note maximale possible} & \textbf{20} \\
\hline

\end{longtable}
}

\end{document}
